% PAAMS 2026 Demonstration Paper
% Template: LNCS/LNAI (Springer)
% Max 6 pages

\documentclass[runningheads]{llncs}
\usepackage{tikz}
\usetikzlibrary{positioning, arrows.meta}
\usepackage[T1]{fontenc}
\usepackage{graphicx}
\usepackage{hyperref}
\usepackage{color}
\usepackage{booktabs}
\usepackage{amsmath}
\usepackage{url}
\renewcommand\UrlFont{\color{blue}\rmfamily}

\begin{document}

\title{A Multi-Agent Decision-Support System for Dairy Farm Energy Efficiency}

\titlerunning{Multi-Agent Dairy Farm Energy Advisor}

\author{%
  \and
}

\authorrunning{}

\institute{%
}

\maketitle

\begin{abstract}
This paper presents a multi-agent decision-support system that transforms a
raw hourly electricity export from a dairy farm into a structured energy
report and a plain-language advisory note, without requiring any domain
expertise from the operator. A coordinator agent orchestrates six specialised task agents, each
with autonomous trigger conditions and an independent decision policy, covering
data ingestion and cleaning, herd-size estimation, statistical energy
characterisation, zero-shot anomaly detection using the MOMENT-1-large
foundation model, knowledge retrieval via a curated dairy energy knowledge
base, and narrative generation using a locally deployed large language model.
The demonstration shows the full pipeline through an interactive web
application: the operator uploads a CSV file and receives, progressively, a
herd-size estimate, an annotated energy dashboard, an anomaly timeline, and a
farm-specific advisory report. A preliminary evaluation with five agricultural
energy professionals confirmed the reports to be factually grounded and
immediately actionable without specialist training.
\end{abstract}

\keywords{Multi-agent systems \and Energy efficiency \and Dairy farming \and
Foundation models \and Retrieval-augmented generation \and Anomaly detection}

% ============================================================
\section{Introduction}
\label{sec:intro}
% ============================================================

Dairy farming is one of the most energy-intensive segments of agriculture.
Milking equipment, bulk-tank refrigeration, hot-water systems for plant
washing, ventilation, and lighting together consume on the order of 300 to
600 kWh per cow per year~\cite{Upton2015}. With electricity prices in Ireland
and across Europe frequently exceeding 0.30 EUR per kWh, energy represents a
significant and growing share of operational costs. Despite this, most farm
operators receive only a monthly aggregated electricity bill, with no
indication of where the load originates, whether consumption is within a
healthy range for their herd size, or when equipment has malfunctioned.

Agent-based models have been applied to simulate the electricity consumption
patterns of Irish dairy farms, capturing how herd size, milking schedules, and
seasonal conditions interact to produce the farm's overall load
profile~\cite{Khaleghy2026}. This body of work provides both validated
synthetic consumption profiles and per-cow energy benchmarks grounded in real
farm operation. However, the outputs of such models have not yet been
integrated into an end-to-end advisory tool that an operator can use directly.

Existing advisory tools either require on-site sub-metering hardware or depend
on manual interpretation by an energy consultant. Automated approaches have
been explored for anomaly detection in energy systems~\cite{Cook2019} and for
retrieval-augmented generation in knowledge-intensive domains~\cite{Lewis2020},
but the combination of consumption-based herd-size estimation, zero-shot
foundation-model anomaly detection, and local LLM narrative generation has not
been demonstrated in an agricultural advisory context.

This paper presents a system that addresses this gap. A coordinator agent
delegates tasks across six specialised agents, requiring only a single
commodity input, namely an hourly electricity CSV export, to produce actionable
advice a non-expert operator can act on directly.

% ============================================================
\section{Main Purpose}
\label{sec:purpose}
% ============================================================

Given only a farm's hourly electricity consumption time series, the system
estimates herd size from the consumption pattern so energy use can be
contextualised per cow against industry benchmarks, characterises consumption
at hourly, monthly and seasonal granularity to expose the baseline idle load
and peak operating hours, detects anomalous hours without any labelled
training data, and generates a plain-language advisory grounded in a curated
knowledge base and tailored to each farm.

\subsection{Multi-Agent Architecture}

The pipeline is implemented in Python and exposed through a Streamlit web
interface. A coordinator agent receives the uploaded CSV, delegates tasks to
six specialised agents, accumulates their outputs in a shared JSON report
object, and applies fallback strategies on failure (Fig.~\ref{fig:arch}). Each
task agent owns its local reasoning, holds an independent decision policy,
and fires autonomously when its input is ready; because all inter-agent
communication passes through the shared JSON, any agent is asynchronously
replaceable.

\begin{figure}[t]
  \centering
  \includegraphics[width=0.90\textwidth]{fig1.pdf}
  \caption{Coordinator-delegated multi-agent pipeline. Six specialist agents fire autonomously and write to a shared JSON report; the RAG agent grounds the narrative in a curated knowledge base..}
  \label{fig:arch}
\end{figure}

The data agent validates the uploaded file, infers measurement units, fills
short gaps by linear interpolation, and hands a clean hourly series to the
herd-size estimation agent. The herd-size agent derives cow count from the
consumption pattern using an ensemble trained on synthetic profiles from the
agent-based model in~\cite{Khaleghy2026}, and when the farm falls outside the
training distribution it switches autonomously to a consumption-based
benchmark strategy and reports a plausible range.

The energy analysis agent computes hourly, monthly and seasonal statistics,
including the baseline idle load, its share of mean hourly consumption, the
annual idle-load cost, and the peak and off-peak hours; every derived figure
is pre-computed here so the downstream narrative agent never performs
arithmetic. The anomaly detection agent applies MOMENT-1-large in
reconstruction mode~\cite{Goswami2024}, sliding overlapping windows across the
series, deriving its flagging threshold from the farm's empirical residual
distribution, and classifying each flagged hour as a high-consumption or
low-consumption event with no farm-specific training.

The retrieval-augmented generation (RAG) agent maintains a curated dairy
energy knowledge base of per-cow benchmarks, baseline-load health ranges,
efficiency measures, anomaly heuristics and seasonal load drivers; it embeds
each section's data payload with all-MiniLM-L6-v2~\cite{Wang2020} and
retrieves the most relevant passages for the LLM prompt. The LLM narrative
agent generates the advisory section by section using a locally deployed
qwen2.5:14b model via Ollama~\cite{Qwen2024}; each section receives only the
fields relevant to it, and strict prompt constraints force the model to
quote pre-computed figures verbatim and cite only benchmarks from the
retrieved guidance. The pipeline runs entirely on CPU.

\subsection{Multi-Agent Design Properties}

The architecture satisfies the defining properties of a multi-agent system.
Each task agent encapsulates independent local reasoning, the anomaly agent
derives its flagging threshold from the farm's own residuals and the
herd-estimation agent switches strategy autonomously when its training
distribution assumption is violated. Each agent has an autonomous trigger
condition tied to data readiness, and loose coupling through the shared JSON
object guarantees that any agent can be swapped for a higher-capability
successor without modifying the rest of the pipeline.

% ============================================================
\section{Demonstration}
\label{sec:demo}
% ============================================================

The demonstration runs through an interactive Streamlit web
application.\footnote{Source code and cached demonstration data:
\url{https://github.com/HosseinKhaleghy/LLM-Report-Generation}}
Attendees may upload their own real hourly electricity CSV file or use a
preloaded dataset representing a real Irish dairy farm. The coordinator
initiates the pipeline immediately, and observers can follow each agent's
output as it arrives, inspect intermediate results across tabs, and download
the advisory report. The goal is to show how a multi-agent system decomposes
a task that no single model handles reliably into a sequence of focused,
auditable agent outputs.

\subsection{Interface and Interaction Flow}

The interface presents six tabs that populate as each task agent completes.
The Overview tab shows the herd-size estimate with plausible range and KPI
cards for total consumption, baseline share, and estimated annual cost. The
Hourly, Monthly and Seasonal tabs visualise consumption at their respective
granularities, with the highest and lowest periods highlighted. The Anomalies
tab shows a timeline of flagged hours coloured by direction, a table of the
largest deviations, and a plain-language severity assessment. The Advisory
tab renders the LLM narrative as Markdown with a one-click download button.

The herd-size estimate and full energy dashboard appear within seconds of
upload, as model weights are pre-loaded at application startup. Anomaly
detection completes in approximately 60 to 120 seconds on CPU, after which
the Anomalies tab becomes interactive. LLM generation takes approximately 3
to 5 minutes on CPU and renders section by section as each completes, giving
attendees time to explore the energy and anomaly tabs while the advisory
produces. A preloaded farm example ensures the demonstration can begin
immediately without any upload delay.

\subsection{Example Output}

Figure~\ref{fig:screenshot} shows the application running on a real Irish
dairy farm dataset with 8521 hours of recorded consumption. The pipeline
estimated a herd of 164 cows (plausible range 115 to 192), a total annual
consumption of 57505.9 kWh (17251.78 EUR at the farm's tariff), and a
baseline idle load of 4679.4 Wh representing 69.3\% of mean hourly
consumption, at an annual idle-load cost of 12297.46 EUR. The anomaly
detection agent flagged 43 hours (0.5\%), with 16 classified as
high-consumption and 27 as low-consumption events.

\begin{figure}[t]
  \centering
  \includegraphics[width=\textwidth]{overview.png}
  \caption{The application running on a real dairy farm dataset. The Overview
    tab (shown) displays the herd-size estimate, KPI cards for consumption
    and cost, and the hour-of-day load profile. The remaining tabs provide
    monthly and seasonal breakdowns, an anomaly timeline, and the full LLM
    advisory.}
  \label{fig:screenshot}
\end{figure}

\subsection{Preliminary Usability Evaluation}

A preliminary evaluation was conducted with five agricultural energy
professionals, each interacting with the system on an anonymised real farm
dataset. All five confirmed that the advisory contained no fabricated figures
and that every recommendation was grounded in a data value verifiable in the
dashboard; four out of five rated the herd-size estimate as plausible; and
all five identified at least one maintenance action from the anomaly timeline
they would prioritise. Participants engaged with the energy and anomaly tabs
while the advisory generated, and none reported the wait as excessive.

To complement the human evaluation, an automated LLM-as-judge protocol was
applied to reports generated for five real dairy farms (farms 2, 4, 5, 6,
and 7 from the dataset). A locally deployed qwen2.5:14b judge model scored
each of the eight report sections on three dimensions: factual grounding
(D1, whether every cited figure is traceable to the farm data), benchmark
integrity (D2, whether all cited thresholds are attributable to the
knowledge base), and physical soundness (D3, whether causal reasoning
matches Irish dairy farm reality). Across all 40 section-level judgements
(five farms, eight sections each), mean scores were D1\,=\,5.00,
D2\,=\,4.63, D3\,=\,4.85, and overall\,=\,4.70 on a 1--5 scale.
Factual grounding was perfect across all farms; the modest D2 variance
reflects isolated instances where the model cited a benchmark range that
the judge could not directly map to a retrieved knowledge-base passage.
Because the narrative generator and the judge are drawn from the same
qwen2.5 model family, these figures should be interpreted as an
optimistic upper bound, and an independent judge from a different model
family would likely score the physical-soundness dimension more
critically.

% ============================================================
\section{Conclusions}
\label{sec:conclusions}
% ============================================================

This paper has demonstrated a multi-agent pipeline in which a coordinator
agent delegates tasks across six specialised agents, each with autonomous
trigger conditions, local reasoning, and an independent decision policy, to
transform a single hourly electricity CSV into a farm-specific energy
advisory. The
system operates without labelled anomaly data, without domain expertise from
the operator, and without cloud GPU resources. Three design choices proved
critical for output quality: pre-computing all derived figures in the energy
analysis agent so that the LLM narrative agent never performs arithmetic;
narrowing the JSON payload for each LLM section to only the relevant fields;
and grounding the narrative in a curated knowledge base through the RAG agent.
Together these choices substantially reduce hallucination and improve factual
fidelity in the advisory text.

A preliminary evaluation with five domain experts confirmed the practical
utility of the output, and an automated LLM-as-judge assessment across five
real farms yielded a mean overall score of 4.70 out of 5, with perfect
factual grounding across all farms. Future work will explore fine-tuning the
narrative agent on expert-validated reports, adding sub-metering
disaggregation, and extending anomaly classification to equipment-specific
fault signatures.

% ============================================================
% References
% ============================================================

\begin{thebibliography}{8}

\bibitem{Upton2015}
Upton, J., Murphy, M., De Boer, I.J.M., Groot Koerkamp, P.W.G., Berentsen,
P.B.M., Shalloo, L.: Investment appraisal of technology innovations on dairy
farm electricity consumption. Journal of Dairy Science 98(2), 898--909 (2015)

\bibitem{Khaleghy2026}
Khaleghy, H., Clifford, E., Mason, K.: An agent based approach to Irish dairy
farm electricity consumption modeling. Renewable Energy Focus, 100866 (2026)

\bibitem{Cook2019}
Cook, A.A., Misirkol, G., Fan, Z., Rana, R.: Anomaly detection for IoT
time-series data: a survey. IEEE Internet of Things Journal 7(7),
6481--6494 (2019)

\bibitem{Lewis2020}
Lewis, P., Perez, E., Piktus, A., et al.: Retrieval-augmented generation for
knowledge-intensive NLP tasks. In: Advances in Neural Information Processing
Systems (NeurIPS), vol.~33, pp.~9459--9474 (2020)

\bibitem{Goswami2024}
Goswami, M., Szafer, K., Choudhry, A., Cai, Y., Li, S., Dubrawski, A.:
MOMENT: a family of open time-series foundation models. In: Proceedings of
the 41st International Conference on Machine Learning (ICML),
pp.~16115--16152 (2024)

\bibitem{Wang2020}
Wang, W., Wei, F., Dong, L., Bao, H., Yang, N., Zhou, M.: MiniLM: deep
self-attention distillation for task-agnostic compression of pre-trained
transformers. In: Advances in Neural Information Processing Systems (NeurIPS),
vol.~33, pp.~5776--5788 (2020)

\bibitem{Qwen2024}
Yang, A., Yang, B., Zhang, B., et al.: Qwen2.5 technical report.
arXiv preprint arXiv:2412.15115 (2024)

\end{thebibliography}

\end{document}
